% GENERATED FILE. Edit canonical lessons, then run npm run generate:books.
% canonical-source-hash: fnv1a64:2cdd8d75d711beb0
% canonical-lessons: RU-C516-melodiya, RU-C516-ritm, RU-C516-pripev, RU-C516-napevat, RU-C516-falshivit, RU-R516-review-509-516

\chapter{Tune and Rhythm}
\label{ch:ru-tune-and-rhythm}
\hlchaptermodality{516}

\begin{chapteropening}
\textbf{By the end of this chapter:} \emph{I can talk about tunes, rhythm and singing in tune.}

Complete the last lesson of chapter 516: I can talk about tunes, rhythm and singing in tune.
\end{chapteropening}

\section[\texorpdfstring{\ru{мелодия} (melódiya)}{мелодия (melódiya)}]{\ru{мелодия} (melódiya) — a melody, a tune}

\label{lesson:RU-C516-melodiya}



\begin{quote}
Before the new one: say the Russian for a string, then the Russian for to tune.
\end{quote}

\subsection*{You'll want to know: \ru{мелодия}}
\textbf{\ru{мелодия}} (\emph{melódiya}) — \textquotedblleft{}a melody, a tune\textquotedblright{}.

\begin{quote}
\ru{Эта} \ru{мелодия} \ru{мне} \ru{очень} \ru{знакома}.
\end{quote}

\begin{quote}
\emph{Éta melódiya mne óchen znakóma.}
\end{quote}

\begin{quote}
This tune is very familiar to me.
\end{quote}

Feminine.

\begin{cousinweb}
\emph{\ru{мелодия}} comes from Greek \emph{melōidía} 'singing, chanting', built on \emph{mélos} 'song' (EtymDB).
\end{cousinweb}

\begin{grammarlens}[title={What you've built}]
One more word for telling, explaining and describing at B1.
\end{grammarlens}

\subsection*{Your turn}
\begin{itemize}
\raggedright
  \item \emph{Say it:} \emph{melódiya}
  \item \emph{Say it:} \emph{melódiya}, once more
  \item \emph{Say it:} a sentence of your own with \emph{melódiya}
  \item \emph{Recall:} say the Russian for a string, then the Russian for to tune, then say \emph{melódiya} again
\end{itemize}

\begin{tcolorbox}[breakable,colback=teal!4,colframe=teal!35!black,title={Before you move on}]
What does \ru{мелодия} mean? (\textquotedblleft{}A melody, a tune\textquotedblright{}.) Say it once more.
\end{tcolorbox}
\section[\texorpdfstring{\ru{ритм} (ritm)}{ритм (ritm)}]{\ru{ритм} (ritm) — rhythm}

\label{lesson:RU-C516-ritm}



\begin{quote}
Before the new one: say the Russian for to tune, then the Russian for a melody.
\end{quote}

\subsection*{You'll want to know: \ru{ритм}}
\textbf{\ru{ритм}} (\emph{ritm}) — \textquotedblleft{}rhythm\textquotedblright{}.

\begin{quote}
\ru{Танцоры} \ru{двигались} \ru{в} \ru{ритме} \ru{музыки}.
\end{quote}

\begin{quote}
\emph{Tantsóry dvígalis v rítme múzyki.}
\end{quote}

\begin{quote}
The dancers moved to the rhythm of the music.
\end{quote}

Masculine.

\begin{cousinweb}
\emph{\ru{ритм}} was borrowed from Latin \emph{rhythmus}, from Greek \emph{rhythmós} 'measured flow', built on \emph{rhéō} 'to flow' (EtymDB).
\end{cousinweb}

\begin{grammarlens}[title={What you've built}]
One more word for telling, explaining and describing at B1.
\end{grammarlens}

\subsection*{Your turn}
\begin{itemize}
\raggedright
  \item \emph{Say it:} \emph{ritm}
  \item \emph{Say it:} \emph{ritm}, once more
  \item \emph{Say it:} a sentence of your own with \emph{ritm}
  \item \emph{Recall:} say the Russian for to tune, then the Russian for a melody, then say \emph{ritm} again
\end{itemize}

\begin{tcolorbox}[breakable,colback=teal!4,colframe=teal!35!black,title={Before you move on}]
What does \ru{ритм} mean? (\textquotedblleft{}Rhythm\textquotedblright{}.) Say it once more.
\end{tcolorbox}
\section[\texorpdfstring{\ru{припев} (pripév)}{припев (pripév)}]{\ru{припев} (pripév) — a chorus, a refrain}

\label{lesson:RU-C516-pripev}



\begin{quote}
Before the new one: say the Russian for a melody, then the Russian for rhythm.
\end{quote}

\subsection*{You'll want to know: \ru{припев}}
\textbf{\ru{припев}} (\emph{pripév}) — \textquotedblleft{}a chorus, a refrain\textquotedblright{}.

\begin{quote}
\ru{Все} \ru{в} \ru{зале} \ru{пели} \ru{припев} \ru{вместе} \ru{с} \ru{певцом}.
\end{quote}

\begin{quote}
\emph{Vse v zále péli pripév vméste s pevtsóm.}
\end{quote}

\begin{quote}
Everyone in the hall sang the chorus with the singer.
\end{quote}

Masculine. The verse is \emph{\ru{куплет}}.

\begin{cousinweb}
\emph{\ru{припев}} is built on \emph{\ru{при}-} and the root of \emph{\ru{петь}} 'to sing'; EtymDB has no entry.
\end{cousinweb}

\begin{grammarlens}[title={What you've built}]
One more word for telling, explaining and describing at B1.
\end{grammarlens}

\subsection*{Your turn}
\begin{itemize}
\raggedright
  \item \emph{Say it:} \emph{pripév}
  \item \emph{Say it:} \emph{pripév}, once more
  \item \emph{Say it:} a sentence of your own with \emph{pripév}
  \item \emph{Recall:} say the Russian for a melody, then the Russian for rhythm, then say \emph{pripév} again
\end{itemize}

\begin{tcolorbox}[breakable,colback=teal!4,colframe=teal!35!black,title={Before you move on}]
What does \ru{припев} mean? (\textquotedblleft{}A chorus, a refrain\textquotedblright{}.) Say it once more.
\end{tcolorbox}
\section[\texorpdfstring{\ru{напевать} (napevát)}{напевать (napevát)}]{\ru{напевать} (napevát) — to hum}

\label{lesson:RU-C516-napevat}



\begin{quote}
Before the new one: say the Russian for rhythm, then the Russian for a chorus.
\end{quote}

\subsection*{You'll want to know: \ru{напевать}}
\textbf{\ru{напевать}} (\emph{napevát}) — \textquotedblleft{}to hum\textquotedblright{}.

\begin{quote}
\ru{На} \ru{кухне} \ru{мама} \ru{напевает} \ru{любимую} \ru{песню}.
\end{quote}

\begin{quote}
\emph{Na kúkhne máma napeváyet lyubímuyu pésnyu.}
\end{quote}

\begin{quote}
Mum is humming her favourite song in the kitchen.
\end{quote}

Imperfective; present \ru{напеваю}. Perfective \ru{напеть}.

\begin{cousinweb}
\emph{\ru{напевать}} is built on \emph{\ru{на}-} and \emph{\ru{петь}} 'to sing'; EtymDB has no entry.
\end{cousinweb}

\begin{grammarlens}[title={What you've built}]
One more word for telling, explaining and describing at B1.
\end{grammarlens}

\subsection*{Your turn}
\begin{itemize}
\raggedright
  \item \emph{Say it:} \emph{napevát}
  \item \emph{Say it:} \emph{napevát}, once more
  \item \emph{Say it:} a sentence of your own with \emph{napevát}
  \item \emph{Recall:} say the Russian for rhythm, then the Russian for a chorus, then say \emph{napevát} again
\end{itemize}

\begin{tcolorbox}[breakable,colback=teal!4,colframe=teal!35!black,title={Before you move on}]
What does \ru{напевать} mean? (\textquotedblleft{}To hum\textquotedblright{}.) Say it once more.
\end{tcolorbox}
\section[\texorpdfstring{\ru{фальшивить} (falshívit)}{фальшивить (falshívit)}]{\ru{фальшивить} (falshívit) — to sing out of tune}

\label{lesson:RU-C516-falshivit}



\begin{quote}
Before the new one: say the Russian for a chorus, then the Russian for to hum.
\end{quote}

\subsection*{You'll want to know: \ru{фальшивить}}
\textbf{\ru{фальшивить}} (\emph{falshívit}) — \textquotedblleft{}to sing out of tune\textquotedblright{}.

\begin{quote}
\ru{Он} \ru{громко} \ru{поёт}, \ru{но} \ru{немного} \ru{фальшивит}.
\end{quote}

\begin{quote}
\emph{On grómko poyót, no nemnógo falshívit.}
\end{quote}

\begin{quote}
He sings loudly but a little out of tune.
\end{quote}

Imperfective; present \ru{фальшивлю}, \ru{фальшивишь} (\ru{в} becomes \ru{вл}).

\begin{cousinweb}
\emph{\ru{фальшивить}} is derived from \emph{\ru{фальшь}} 'falseness' (EtymDB), which was borrowed from German \emph{falsch} 'false'.
\end{cousinweb}

\begin{grammarlens}[title={What you've built}]
One more word for telling, explaining and describing at B1.
\end{grammarlens}

\subsection*{Your turn}
\begin{itemize}
\raggedright
  \item \emph{Say it:} \emph{falshívit}
  \item \emph{Say it:} \emph{falshívit}, once more
  \item \emph{Say it:} a sentence of your own with \emph{falshívit}
  \item \emph{Recall:} say the Russian for a chorus, then the Russian for to hum, then say \emph{falshívit} again
\end{itemize}

\begin{tcolorbox}[breakable,colback=teal!4,colframe=teal!35!black,title={Before you move on}]
What does \ru{фальшивить} mean? (\textquotedblleft{}To sing out of tune\textquotedblright{}.) Say it once more.
\end{tcolorbox}
\section[(dialogue)]{Words from chapters 509–516 — a review}

\label{lesson:RU-R516-review-509-516}



\begin{quote}
Say the Russian for gates and a goalkeeper.
\end{quote}

\subsection*{Your turn}
Say these aloud:

\begin{itemize}
\raggedright
  \item gates, a goalkeeper, a forward, a goal, to score
  \item a defeat, a tie, a rival, a loss, to be in the lead
  \item a tournament, a championship, a champion, a medal, to award
  \item a marathon, a runner, the finish, to fall behind, a distance
  \item a warm-up, to do press-ups, to squat, dumbbells, stamina
  \item a skating rink, a figure skater, to slide, flippers, to splash about
  \item a grand piano, a cello, a flute, a string, to tune
  \item a melody, rhythm, a chorus, to hum, to sing out of tune
\end{itemize}

\begin{tcolorbox}[breakable,colback=teal!4,colframe=teal!35!black,title={Before you move on}]
Gates? (\textbf{\ru{ворота}}.) A goalkeeper? (\textbf{\ru{вратарь}}.) A forward? (\textbf{\ru{нападающий}}.) A goal? (\textbf{\ru{гол}}.) To score? (\textbf{\ru{забивать}}.) A defeat? (\textbf{\ru{поражение}}.) A tie? (\textbf{\ru{ничья}}.) A rival? (\textbf{\ru{соперник}}.) A loss? (\textbf{\ru{проигрыш}}.) To be in the lead? (\textbf{\ru{лидировать}}.) A tournament? (\textbf{\ru{турнир}}.) A championship? (\textbf{\ru{чемпионат}}.) A champion? (\textbf{\ru{чемпион}}.) A medal? (\textbf{\ru{медаль}}.) To award? (\textbf{\ru{награждать}}.) A marathon? (\textbf{\ru{марафон}}.) A runner? (\textbf{\ru{бегун}}.) The finish? (\textbf{\ru{финиш}}.) To fall behind? (\textbf{\ru{отставать}}.) A distance? (\textbf{\ru{дистанция}}.) A warm-up? (\textbf{\ru{разминка}}.) To do press-ups? (\textbf{\ru{отжиматься}}.) To squat? (\textbf{\ru{приседать}}.) Dumbbells? (\textbf{\ru{гантели}}.) Stamina? (\textbf{\ru{выносливость}}.) A skating rink? (\textbf{\ru{каток}}.) A figure skater? (\textbf{\ru{фигурист}}.) To slide? (\textbf{\ru{скользить}}.) Flippers? (\textbf{\ru{ласты}}.) To splash about? (\textbf{\ru{плескаться}}.) A grand piano? (\textbf{\ru{рояль}}.) A cello? (\textbf{\ru{виолончель}}.) A flute? (\textbf{\ru{флейта}}.) A string? (\textbf{\ru{струна}}.) To tune? (\textbf{\ru{настраивать}}.) A melody? (\textbf{\ru{мелодия}}.) Rhythm? (\textbf{\ru{ритм}}.) A chorus? (\textbf{\ru{припев}}.) To hum? (\textbf{\ru{напевать}}.) To sing out of tune? (\textbf{\ru{фальшивить}}.)
\end{tcolorbox}
